% ============================================================================
\clearpage
\section*{v4.9.0-r14 Source-Only: Why Six Kicks per Beat --- Readout Directions from the Triole and the Umkehrwalze}
\addcontentsline{toc}{section}{v4.9.0-r14 Source-Only: Why Six Kicks per Beat --- Readout Directions from the Triole and the Umkehrwalze}
\label{sec:v490-r14-readout-directions}
% ============================================================================

\paragraph{Purpose.}
r13 fixed the kick interval of the dynamics of r11 at one sixth of the clock beat.  It left open why the compensator acts
six times per beat.  This layer offers a reason built from two objects of the source: the triole and the involutive
Umkehrwalze \(J_B\) of r02.3.  The reason is an interpretation (admitted by the author, 2026-10-22), not a derivation.  The
script is \texttt{llm/k55\_kick\_richtungen.py}.  Earlier layers are not rewritten.

\noindent\textbf{Source-only status.}
This v4.9.0-r14 layer is source-only.  It does not build a PDF, does not assign a new DOI,
and does not authorize a public release.  The v4.9.0 Version DOI
\href{https://doi.org/10.5281/zenodo.23534954}{10.5281/zenodo.23534954} and the Concept DOI
\href{https://doi.org/10.5281/zenodo.19210728}{10.5281/zenodo.19210728} are unchanged by it.

\subsection*{v4.9.0-r14 kicks as passages of readout directions}
\label{subsec:v490-r14-passages}

Replace the periodic kick of r11 by an event: the compensator acts whenever the phase of the invariant \(s\) passes one of a
fixed set of readout directions.  Under a uniform duality rotation the events are periodic, and their spacing is set by the
set of directions.

\paragraph{Theoremlet v4.9.0-r14.1.}
In the \(s\)-plane, \(\arg s=2\alpha\) with \(\alpha\) the duality angle.
\begin{enumerate}[leftmargin=2em]
\item[(a)] The triole directions \(\arg s\in\{0^\circ,120^\circ,240^\circ\}\) and their images under \(J_B\), which acts as
\(s\mapsto-\bar s\) (r02.3), are together the six directions \(0^\circ,60^\circ,\dots,300^\circ\).  The set contains matter
(\(0^\circ\)) and tachyon (\(180^\circ\)).
\item[(b)] Passing them under a uniform rotation gives kicks every \(30^\circ\) of duality angle, hence the stage factor
\(\cos^2 30^\circ=3/4\).  The triole alone gives \(60^\circ\) and \(q=1/4\).  A quarter grid gives \(45^\circ\) and \(q=1/2\).
Only triole and \(J_B\) together give the factor of the source.
\item[(c)] One revolution of \(s\) around the photon (\(\arg s\) from \(0^\circ\) to \(360^\circ\)) passes the six directions once.
It is the duality half turn \(z\mapsto-z\), the charge conjugation of r02.2, and by r13 it lasts exactly one beat.
\end{enumerate}

\paragraph{Proof sketch.}
\(-\overline{e^{i\theta}}=e^{i(\pi-\theta)}\) maps \(\{0,120,240\}\) to \(\{180,60,300\}\).  The rest follows from
\(\arg s=2\alpha\) and r10.1.  Checked numerically.  \(\square\)

\paragraph{Reading.}
The six kicks per beat are then not an external schedule but the six readout directions met by one revolution of \(s\)
around the photon.  The factor \(3/4\), the detector stage \(20\), the beat, the quarter turn from matter to tachyon and the
charge conjugation of the round trip become aspects of one construction: a uniform duality rotation read out at
\(\text{triole}\cup J_B(\text{triole})\).

\paragraph{What is assumed.}
The triole of the source (Theoremlet r06.2) is a condition on the three components of a vector, \((1,\omega,\omega^2)\), not on
the phase of \(s\).  Placing the triole on \(\arg s\) is the interpretive step of this layer.  Squaring maps the triole of
phases onto a triole of phases, so the placement is consistent, but the source does not make it.  Why the readout should
act at discrete directions at all is also assumed.

\paragraph{Remark, offered and not adopted.}
Six directions at \(60^\circ\) with the triole permutations as symmetry are the root system \(A_2\) of \(SU(3)\), whose Weyl
group is \(S_3\).  r15 of v4.8.0 met \(SU(3)\) as the stabiliser of the symmetric point, there read as colour.  Whether the
readout hexagon and the colour \(SU(3)\) are related would need a map from the duality plane to a Cartan plane, which the
source does not supply.

\subsection*{v4.9.0-r14 register entries}
\label{subsec:v490-r14-register}

\begin{center}
\small
\begin{tabular}{@{}>{\raggedright\arraybackslash}p{0.34\textwidth}>{\raggedright\arraybackslash}p{0.58\textwidth}@{}}
\toprule
statement & status after v4.9.0-r14\\
\midrule
six kicks per beat (r13) & reading offered: readout at \(\text{triole}\cup J_B(\text{triole})\) in the \(s\)-plane (r14.1)\\
factor \(3/4\) & needs both triole and \(J_B\); either alone, or a quarter grid, gives another factor\\
beat & one revolution of \(s\) around the photon, equal to the charge conjugation of r02.2\\
\(A_2\)/colour & remark only\\
\bottomrule
\end{tabular}
\end{center}

\subsection*{v4.9.0-r14 ledger and guard update}
\label{subsec:v490-r14-ledger-guard}

\[
R_{\rm kick}^{v490\text{-}r14}=(0,0,1)
\]
for (a reason for the kick schedule offered, its consequences consistent with r10--r13, the placement of the triole open).

The layer spends no guard:
\[
R_{\rm glob}^{(v490\text{-}r14,+)}=(0,1,1,1,1),
\qquad
\Gamma_{\rm guard}^{v490\text{-}r14}=(1,1,1,1,1),
\qquad
\beta_4=1/256\ \text{unspent}.
\]
It does not imply exact finite \(C5\) zero closure, exact all-readout terminality,
\(C6/T2\), a global photon-clock/RH theorem, or global all-readout closure.

